\chapter{Dynamic Programming: DP on Squares in a Matrix}
\chapterepigraph{This closing chapter asks one last question about a
binary matrix: how large a square (or rectangle) of all ones can be
found, anchored at each cell? Once a cell knows the best answer for
its neighbors above, left, and diagonally above-left, it can answer for
itself in one comparison -- no search over a range required.}

\section{The Three-Neighbor Square DP}
\label{sec:dp-squares-intro}

The defining recurrence of this chapter looks at exactly three already-computed
neighbors -- directly above, directly left, and diagonally above-left --
and takes their minimum: a square of all ones can only be extended by
one row and one column simultaneously if \emph{all three} of those
smaller squares already existed.

\begin{patternbox}[The Two Problems in This Chapter]
\begin{itemize}
  \item \textbf{Count Square Submatrices with All Ones (Section~\ref{sec:dp-count-square-submatrices})}:
        the three-neighbor-minimum recurrence directly, summed over
        every cell.
  \item \textbf{Maximal Rectangle (Section~\ref{sec:dp-maximal-rectangle})}:
        a related but distinct DP -- computing, for each cell, the
        height of consecutive ones \emph{above} it (a simple 1D
        recurrence per column), which then feeds into the largest-
        rectangle-in-a-histogram technique already developed in the
        Stack and Queue chapter.
\end{itemize}
\end{patternbox}

\newpage

\section{Count Square Submatrices with All Ones}
\label{sec:dp-count-square-submatrices}

\problemstatement{Given a binary matrix $\textit{mat}$, count the total
number of square submatrices that contain only $1$s (squares of every
size, at every position).}

\begin{patternbox}
\emph{``Largest/count of square submatrices of all ones, anchored at
each cell''} is solved by a recurrence that looks at exactly three
neighbors: a square of side $s$ ending at $(i,j)$ as its bottom-right
corner exists only if squares of side $s-1$ already exist ending at
\textbf{all three} of the cell above, the cell to the left, and the
cell diagonally above-left.
\end{patternbox}

\subsection*{Deriving the recurrence}

Let $\textit{dp}[i][j]$ be the side length of the \textbf{largest}
all-ones square whose bottom-right corner is exactly $(i,j)$ (and
$0$ if $\textit{mat}[i][j]=0$). If $\textit{mat}[i][j]=1$, extending a
square one row down and one column right requires the square
\emph{ending} at each of the three neighboring corners to already
support at least side length $\textit{dp}[i][j]-1$:
\[
  \textit{dp}[i][j] = \begin{cases}
  0 & \textit{mat}[i][j]=0, \\
  1+\min\big(\textit{dp}[i{-}1][j],\ \textit{dp}[i][j{-}1],\
  \textit{dp}[i{-}1][j{-}1]\big) & \textit{mat}[i][j]=1.
  \end{cases}
\]

\begin{ideabox}[Why the Minimum, and Why Summing Gives the Count]
The new square's side is capped by whichever of the three neighbors
supports the \emph{smallest} square -- a single weak corner limits how
far the square can extend, exactly like the weakest link in a chain.
Separately, a cell with $\textit{dp}[i][j]=k$ doesn't just mean
\emph{one} square ends there -- it means $k$ different squares end
there: one each of side $1, 2, \ldots, k$ (every smaller square nested
inside the largest one is also a valid, distinct all-ones square with
its bottom-right corner at $(i,j)$). So the \textbf{count} of all
square submatrices is simply $\sum_{i,j} \textit{dp}[i][j]$.
\end{ideabox}

\subsection*{C++ implementation}

\begin{lstlisting}
#include <bits/stdc++.h>
using namespace std;

int countSquareSubmatrices(vector<vector<int>>& mat) {
    int n = mat.size(), m = mat[0].size();
    vector<vector<int>> dp(n, vector<int>(m, 0));
    int totalCount = 0;

    for (int i = 0; i < n; i++) {
        for (int j = 0; j < m; j++) {
            if (mat[i][j] == 0) {
                dp[i][j] = 0;
            } else if (i == 0 || j == 0) {
                dp[i][j] = 1; // first row/column: at most a 1x1 square
            } else {
                dp[i][j] = 1 + min({dp[i - 1][j], dp[i][j - 1], dp[i - 1][j - 1]});
            }
            totalCount += dp[i][j];
        }
    }
    return totalCount;
}

int main() {
    vector<vector<int>> mat = {
        {0, 1, 1, 1},
        {1, 1, 1, 1},
        {0, 1, 1, 1}
    };
    cout << countSquareSubmatrices(mat) << endl; // 15
    return 0;
}
\end{lstlisting}

\subsection*{Dry run}

\dryrun{Take the $3\times4$ matrix above.}

\begin{center}
\begin{tabular}{c|c|c|c|c}
 & 0 & 1 & 2 & 3 \\ \hline
0 & 0 & 1 & 1 & 1 \\ \hline
1 & 1 & 1 & 2 & 2 \\ \hline
2 & 0 & 1 & 2 & 3
\end{tabular}
\end{center}

$\textit{dp}[1][2]=2$: $\textit{mat}[1][2]=1$, and
$\min(\textit{dp}[0][2],\textit{dp}[1][1],\textit{dp}[0][1])=
\min(1,1,1)=1$, so $\textit{dp}[1][2]=2$ (a $2\times2$ all-ones square
exists with corners at rows $0$--$1$, columns $1$--$2$).
$\textit{dp}[2][3]=3$: neighbors $\textit{dp}[1][3]=2$,
$\textit{dp}[2][2]=2$, $\textit{dp}[1][2]=2$, minimum $2$, giving
$\textit{dp}[2][3]=3$ (a full $3\times3$ all-ones square, columns
$1$--$3$, all three rows). Summing every cell:
$0{+}1{+}1{+}1{+}1{+}1{+}2{+}2{+}0{+}1{+}2{+}3 = 15$.

\begin{complexitybox}
\textbf{Time Complexity.} $O(n \cdot m)$ -- one constant-time
comparison per cell.
\textbf{Space Complexity.} $O(n \cdot m)$ for the \textit{dp} table
(reducible to $O(m)$ with a rolling pair of rows, since each cell only
reads the row directly above and the current row).
\end{complexitybox}

\begin{takeawaybox}
A surprisingly large class of ``2D shape ending here'' problems reduce
to a three-neighbor minimum (or maximum) plus one: the moment a shape
can only grow when \emph{all} of its supporting smaller shapes already
exist, the recurrence is this same one-line comparison, with no search
over a range needed at all -- the simplest recurrence shape in this
entire book.
\end{takeawaybox}

\section{Maximal Rectangle}
\label{sec:dp-maximal-rectangle}

\problemstatement{Given a binary matrix $\textit{mat}$, find the area
of the largest rectangle (not necessarily square) containing only
$1$s.}

\begin{patternbox}
This closing problem of the entire book splits cleanly into two
independent techniques: a simple per-column \textbf{height DP} (how
many consecutive $1$s sit directly above and including this cell),
followed by treating each row's height array as a \textbf{histogram}
and finding its largest rectangle -- exactly the monotonic-stack
technique already developed for Largest Rectangle in Histogram in the
Stack and Queue chapter.
\end{patternbox}

\subsection*{Deriving the height DP}

Let $\textit{height}[i][j]$ be the number of consecutive $1$s ending at
row $i$, looking upward, in column $j$:
\[
  \textit{height}[i][j] = \begin{cases}
  0 & \textit{mat}[i][j]=0, \\
  \textit{height}[i{-}1][j] + 1 & \textit{mat}[i][j]=1.
  \end{cases}
\]
This is the same ``reset to $0$ on failure, otherwise extend by $1$''
shape seen throughout this book (Longest Common Substring's mismatch
reset, House Robber's running totals) -- here applied independently to
every column.

\subsection*{Reframing each row as a histogram}

After computing $\textit{height}[i][\cdot]$ for a fixed row $i$, that
row \emph{is} a histogram: $\textit{height}[i][j]$ is the bar height at
position $j$. The largest all-ones rectangle with its \textbf{bottom
edge on row $i$} is exactly the largest rectangle in that histogram --
solved in $O(m)$ using the monotonic-stack technique from the Stack and
Queue chapter's Largest Rectangle in Histogram problem. Running this
once per row and tracking the global maximum solves the full problem.

\subsection*{C++ implementation}

\begin{lstlisting}
#include <bits/stdc++.h>
using namespace std;

int largestRectangleInHistogram(vector<int>& heights) {
    int n = heights.size();
    stack<int> st;
    int maxArea = 0;

    for (int i = 0; i <= n; i++) {
        int currentHeight = (i == n) ? 0 : heights[i];
        while (!st.empty() && heights[st.top()] >= currentHeight) {
            int height = heights[st.top()];
            st.pop();
            int width = st.empty() ? i : i - st.top() - 1;
            maxArea = max(maxArea, height * width);
        }
        st.push(i);
    }
    return maxArea;
}

int maximalRectangle(vector<vector<int>>& mat) {
    int n = mat.size(), m = mat[0].size();
    vector<int> height(m, 0);
    int best = 0;

    for (int i = 0; i < n; i++) {
        for (int j = 0; j < m; j++) {
            height[j] = (mat[i][j] == 1) ? height[j] + 1 : 0;
        }
        best = max(best, largestRectangleInHistogram(height));
    }
    return best;
}

int main() {
    vector<vector<int>> mat = {
        {1, 0, 1, 0, 0},
        {1, 0, 1, 1, 1},
        {1, 1, 1, 1, 1},
        {1, 0, 0, 1, 0}
    };
    cout << maximalRectangle(mat) << endl; // 6
    return 0;
}
\end{lstlisting}

\subsection*{Dry run}

\dryrun{Take the $4\times5$ matrix above.}

Row-by-row height arrays: row $0$: $[1,0,1,0,0]$. row $1$:
$[2,0,2,1,1]$. row $2$: $[3,1,3,2,2]$. row $3$: $[4,0,0,3,0]$.

At row $2$, the histogram $[3,1,3,2,2]$ has its largest rectangle as
columns $2$--$4$ at height $2$ (area $6$): bars of height $3,2,2$, the
minimum height $2$ spanning width $3$, giving area $6$. This is the overall maximum: at row $0$ the best is area $1$; at row
$1$, best is $3\times1=3$ (columns $2$--$4$, heights $2,1,1$, limited
by the minimum height $1$); at row $3$, best is $4$ (column $0$ alone,
height $4$). The global maximum across all rows is $6$.

\begin{complexitybox}
\textbf{Time Complexity.} $O(n \cdot m)$ -- $n$ rows, each requiring an
$O(m)$ height update and an $O(m)$ histogram pass (the monotonic stack
visits each bar a bounded number of times).
\textbf{Space Complexity.} $O(m)$ for the \textit{height} array and the
stack.
\end{complexitybox}

\begin{takeawaybox}
This book closes on a problem that is itself a closing argument for
everything in it: recognizing that a 2D problem decomposes into ``a
simple DP recurrence run independently per column'' \emph{composed
with} ``a classical 1D technique from an earlier chapter, run once per
row'' is, in the end, the same skill exercised throughout this entire
book -- noticing which already-solved shape a new problem's structure
matches, rather than deriving a solution from nothing.
\end{takeawaybox}

